% ============================================================
% Section 4: SNR Framework (Medium length for TNNLS)
% Target: ~4 pages
% ============================================================

\section{Signal-to-Noise Ratio Framework}
\label{sec:snr}

We develop an information-theoretic framework that explains when graph structure helps or harms node classification, based on the Signal-to-Noise Ratio (SNR) of message passing.

\begin{table}[t]
\centering
\caption{Key Notation Used in This Section}
\label{tab:notation}
\small
\begin{tabular}{@{}cl@{}}
\toprule
\textbf{Symbol} & \textbf{Definition} \\
\midrule
$h$ & Edge homophily, $h = |\{(u,v) \in E : y_u = y_v\}| / |E|$ \\
$\rho$ & Edge correlation coefficient, $\rho = 2h - 1 \in [-1, 1]$ \\
$d$ & Fixed number of aggregated neighbors \\
$s$ & Feature strength, $s=\mu^\top\Sigma^{-1}\mu$ \\
$\kappa$ & Exact discriminability ratio in Proposition~\ref{prop:fixed_degree} \\
$\mathcal{D}(Z)$ & Discriminability of representation $Z$ \\
$\Delta_Z$ & Class separation, $\mathbb{E}[Z|Y=+1] - \mathbb{E}[Z|Y=-1]$ \\
$\Sigma_Z$ & Within-class covariance, $\mathrm{Cov}(Z|Y)$ \\
$\mathcal{B}$ & Information Budget, $\mathcal{B} = 1 - \text{Acc}_{\text{MLP}}$ \\
$\rho_C$ & Effective correlation for $C$ classes, $\rho_C = (Ch-1)/(C-1)$ \\
\bottomrule
\end{tabular}
\end{table}

\subsection{Problem Setup}

\begin{definition}[Sparse Contextual Stochastic Block Model]
\label{def:csbm}
A sparse CSBM graph $G = (V, E)$ with $n$ nodes consists of:
\begin{itemize}
    \item Binary labels $Y_v \in \{-1, +1\}$ with $\mathbb{P}(Y_v = +1) = 1/2$
    \item Features $X_v = Y_v \mu + \varepsilon_v$ where $\varepsilon_v \sim \mathcal{N}(0, \Sigma)$ i.i.d.
    \item Edges: $\mathbb{P}((u,v) \in E | Y_u, Y_v) = \frac{a}{n}$ if $Y_u = Y_v$, else $\frac{b}{n}$
    \item Homophily: $h = a/(a+b)$, average degree: $k = (a+b)/2$
    \item Edge correlation: $\rho = 2h - 1 \in [-1, 1]$
\end{itemize}
\end{definition}

The edge correlation $\rho$ measures the expected label agreement given an edge:
\begin{equation}
    \rho = \mathbb{E}[Y_u | (u,v) \in E, Y_v = +1] = \frac{a-b}{a+b}
\end{equation}
When $h > 0.5$ ($\rho > 0$), neighbors tend to share labels (homophily). When $h < 0.5$ ($\rho < 0$), neighbors tend to have opposite labels (heterophily).

\begin{definition}[Discriminability]
\label{def:discriminability}
The discriminability of representation $Z$ is the squared Mahalanobis distance between class-conditional distributions:
\begin{equation}
    \mathcal{D}(Z) := \Delta_Z^\top \Sigma_Z^{-1} \Delta_Z
\end{equation}
where $\Delta_Z = \mathbb{E}[Z|Y=+1] - \mathbb{E}[Z|Y=-1]$ is the class separation and $\Sigma_Z = \mathrm{Cov}(Z|Y)$ is within-class covariance.
\end{definition}

\begin{proposition}[Bayes Error]
Under Gaussian assumptions, the Bayes-optimal error is:
\begin{equation}
    P_e = \Phi\left(-\frac{1}{2}\sqrt{\mathcal{D}}\right)
\end{equation}
where $\Phi$ is the standard normal CDF. Thus higher $\mathcal{D}$ implies lower error.
\end{proposition}

\subsection{Exact Fixed-Degree Calculation}

We first analyze a deliberately restricted model without replacing a random degree by its expectation. Conditional on a center label $Y_v=y$, let $Y_1,\ldots,Y_d$ be independent neighbor labels satisfying $\mathbb{P}(Y_i=y\mid Y_v=y)=(1+\rho)/2$. Let $X_i=Y_i\mu+\varepsilon_i$, where the $\varepsilon_i\sim\mathcal{N}(0,\Sigma)$ are independent of all labels, and define $A_v=d^{-1}\sum_{i=1}^d X_i$.

\begin{proposition}[Exact Fixed-Degree Discriminability]
\label{prop:fixed_degree}
For $s=\mu^\top\Sigma^{-1}\mu$, the conditional moments and discriminability ratio are
\begin{align}
    \mathbb{E}[A_v\mid Y_v=y] &= \rho y\mu,\\
    \operatorname{Cov}(A_v\mid Y_v=y)
        &= \frac{1}{d}\left(\Sigma+(1-\rho^2)\mu\mu^\top\right),\\
    \kappa(d,\rho,s)
        :=\frac{\mathcal{D}(A)}{\mathcal{D}(X)}
        &=\frac{\rho^2d}{1+(1-\rho^2)s}.
\end{align}
\end{proposition}

\begin{proof}
We derive this result in three steps.

Since $\mathbb{E}[Y_i\mid Y_v=y]=\rho y$ and
$\operatorname{Var}(Y_i\mid Y_v=y)=1-\rho^2$, independence gives the first two identities. Moreover, $\Delta_X=2\mu$, $\mathcal{D}(X)=4s$, and $\Delta_A=2\rho\mu$. Applying the Sherman--Morrison identity,
\begin{equation}
 \mu^\top\left(\Sigma+(1-\rho^2)\mu\mu^\top\right)^{-1}\mu
 =\frac{s}{1+(1-\rho^2)s}.
\end{equation}
Substitution yields the stated ratio.
\end{proof}

\begin{corollary}[Aggregation Threshold]
\label{cor:threshold}
In the fixed-degree model, aggregation improves Mahalanobis discriminability exactly when
\begin{align}
    \kappa>1
    &\iff \rho^2 > \frac{1+s}{d+s}.
\end{align}
Thus the transition depends on feature strength as well as degree and edge correlation. The simpler condition $\rho^2d>1$ is recovered only as $s\to0$.
\end{corollary}

\begin{table}[t]
\centering
\caption{Exact Danger Zone Depends on Feature Strength}
\label{tab:threshold}
\begin{tabular}{@{}ccc@{}}
\toprule
$(d,s)$ & Threshold $|2h-1|$ & Danger Zone ($h$) \\
\midrule
$(10,0)$ & 0.316 & (0.34, 0.66) \\
$(10,1)$ & 0.426 & (0.29, 0.71) \\
$(10,4)$ & 0.598 & (0.20, 0.80) \\
$(20,1)$ & 0.309 & (0.35, 0.65) \\
\bottomrule
\end{tabular}
\end{table}

\textbf{Interpretation} (Table~\ref{tab:threshold}): stronger self-feature separation increases the label-mixture penalty and widens the region in which neighbor-only aggregation loses discriminability. This is a property of the stated moment criterion, not a universal guarantee for a trained GNN.

\subsection{Approximation Regime}

The exact ratio satisfies
\begin{equation}
 \frac{\rho^2d}{1+\eta}\leq\kappa\leq\rho^2d,
 \qquad \eta=(1-\rho^2)s.
\end{equation}
Consequently, replacing $\kappa$ by $\rho^2d$ has relative error exactly $\eta/(1+\eta)$. This makes the required feature-noise-dominance regime explicit and does not invoke an asymptotic $o(1)$ term.

\subsection{Multi-Layer Analysis}

For broadcasting on a regular tree with forward branching factor $b$, the classical Kesten--Stigum quantity is $b\rho^2$ \cite{kesten1966additional,mossel2015reconstruction}. For a $q$-regular undirected tree, $b=q-1$; for a Poisson Galton--Watson limit with mean offspring $k$, $b=k$. We use this only to contextualize the one-hop calculation. We do not claim a multi-layer equality for ordinary GNN aggregation, whose reused neighborhoods, self-loops, nonlinearities, and learned transformations require a separate analysis.

\subsection{Extension to Multi-Class Problems}

For $C$-class problems with uniform within-class and between-class edge probabilities, we propose the following heuristic extension. The effective edge correlation becomes:
\begin{equation}
    \rho_C = \frac{Ch - 1}{C - 1}
\end{equation}
where $h$ is the homophily ratio. The threshold condition becomes $\rho_C^2 k > 1$.

\textbf{Justification.} In the binary case ($C = 2$), we have $\rho_2 = (2h - 1)/(2-1) = 2h - 1 = \rho$, recovering the exact formula. For general $C$, this formula interpolates between the random baseline ($h = 1/C \Rightarrow \rho_C = 0$) and perfect homophily ($h = 1 \Rightarrow \rho_C = 1$). While a rigorous derivation for multi-class CSBM requires tensor analysis beyond the scope of this paper, empirical validation (Section~\ref{sec:experiments}) shows this heuristic accurately predicts GNN behavior on multi-class datasets.

\begin{example}[Cora Dataset]
Cora has $C = 7$ classes and $h = 0.81$. Thus $\rho_7 = (7 \times 0.81 - 1)/(7-1) = 0.78$. With average degree $k \approx 4$, we get $\kappa_7 = 0.78^2 \times 4 = 2.43 > 1$, predicting aggregation is beneficial.
\end{example}

\subsection{Scope and Random-Degree Extension}

The class-conditional distribution of $A_v$ is generally a Gaussian mixture because the neighbor labels are random. Therefore the Gaussian Bayes-error formula stated for $X_v$ must not be applied to $A_v$ solely from its first two moments. Proposition~\ref{prop:fixed_degree} concerns Mahalanobis discriminability only.

For sparse CSBM, the degree is random. Conditional on $D_v=m>0$, Proposition~\ref{prop:fixed_degree} applies with $d=m$. An unconditional calculation must specify the treatment of isolated nodes and average over the degree-conditioned mixture; replacing $D_v$ by its mean is not exact. We leave that random-degree extension separate rather than attach unsupported finite-$n$ error rates to the fixed-degree result.
